    
\documentclass[11pt]{article}
\usepackage{times}
    \usepackage{fullpage}
    
    \title{Aiding diagnosis of Multiple Sclerosis with deep learning}
    \author{Mark McNaught - 2764158M}

    \begin{document}
    \maketitle
    
    
\section*{Project Outline}\label{proposal}

\subsection*{Motivation}\label{motivation}
Multiple Sclerosis (MS) is a complex neurological disorder whose diagnosis can be challenging and 
time-consuming. Deep learning, particularly Convolutional Neural Networks (CNNs), offers the 
opportunity to automate and potentially improve the accuracy and speed of diagnosis from MRI scans, 
which are a primary diagnostic tool. This is especially important for the health sector and for 
patients, as faster and more reliable diagnosis can lead to earlier interventions, improved patient 
outcomes, and reduced clinician workload.



\subsection*{Aims}\label{aims}

The project aims to evaluate and compare deep learning architectures for automatic MS diagnosis 
from MRI scans, determining which model configurations offer the most accurate and robust 
performance. This includes implementing baseline CNNs, transfer-learning-based models, and 
attention-enhanced variants; analysing their training behaviour; and identifying optimised 
architectures. Where feasible, the project may also explore segmentation approaches to isolate 
lesions or relevant regions, providing additional insight into MS diagnosis. Success will be measurable 
through quantitative performance metrics including confusion matrices, accuracy, sensitivity/recall, 
specificity, precision, and F1 scores, alongside clear comparative evaluation across model families 
and evidence of improved diagnostic capability relative to baseline methods.


\section*{Progress}\label{progress}
\begin{itemize}
    % \item Background research conducted on MS, MRI imaging, medical imaging diagnosis procedures, foundational machine learning concepts, and more specific ML approaches used in medical diagnostics.
    % \item Completed a Stanford online machine learning course as well as several introductory ML projects (including a cats vs dogs classifier and a lung cancer detection model) to build the required foundational knowledge.
    % \item Researched Convolutional Neural Networks (CNNs), transfer learning approaches, ImageNet architectures, and their applicability for medical imaging tasks.
    % \item Reviewed research papers on MS diagnosis, medical imaging, and deep learning approaches used in diagnostic pipelines.
    % \item Examined and further researched the provided dataset to determine key data qualities and characteristics.
    % \item Implemented a basic CNN model to assess the complexity of the problem and act as a baseline for future comparisons.
    % \item Implemented several ResNet18 model variants, including attention-augmented architectures using Convolutional Block Attention Modules (CBAMs), and augmented some models with Neighbourhood Components Analysis (NCA) and K-Nearest Neighbours (KNN).
    % \item Tested, analysed, and compared training behaviour, overfitting and underfitting risks, and overall architecture performance.
    % \item Completed initial research into segmentation techniques, RadImageNet for potential medical-specific transfer learning, additional datasets, and applications in medical imaging problems.
    % \item Produced an initial draft of the background section for the dissertation write-up.
    \item \textbf{Background Research and Foundational Concepts:} Completed an online machine learning course, conducted significant background research into Multiple Sclerosis (MS) and its presentation in MRI scans, machine learning in medical diagnostics, CNNs, and transfer learning.
    \item \textbf{Initial Machine Learning Practice:} Completed several introductory ML projects (including a cats vs dogs classifier and a lung cancer detection model) to build the required foundational knowledge.
    \item \textbf{Initial Model Implementation and Experimentation:} Implemented and experimented with a basic CNN model and a ResNet18 model with a custom head layer, reviewing initial performance metrics.
    \item \textbf{Transfer Learning and Model Refinement:} Researched and implemented transfer learning with a ResNet18 model, including modifications such as an early training cut-off for the custom head, and trained the whole model for fewer epochs.
    \item \textbf{Attention Mechanisms:} Implemented a ResNet18 model incorporating Convolutional Block Attention Modules (CBAM), which showed improved results over the basic ResNet18 approach.
    \item \textbf{Advanced Techniques:} Created models implementing Neighborhood Component Analysis (NCA) and k-Nearest Neighbours (kNN) algorithms on top of existing architectures like the CBAM model.
    \item \textbf{Architecture Performance Comparisons:} Reviewed performance across different model architectures to determine most promising approach.
    \item \textbf{Dataset Analysis:} Researched the origin of the project dataset and began investigating potential data augmentation and balancing strategies.
    \item \textbf{Write-up:} Began drafting the initial background section of the write-up.
\end{itemize}

\section*{Problems and Risks}\label{problems-and-risks}

\subsection*{Problems}\label{problems}
\begin{itemize}
    \tightlist
    \item \textbf{Initial Lack of Knowledge and Experience:} No prior experience of machine learning which required significant time spent developing foundational understanding.
    \item \textbf{Computational Limitations:} Experienced limitations with available processing resources and training time, particularly when working with larger models, necessitating the use of smaller variants like ResNet18.
    \item \textbf{Dataset Limitations:} The dataset is relatively small with potential issues from class imbalance (~1400 MS images and ~2100 non-MS images), which can lead to limited generalisability and potential overfitting.
    \item \textbf{Minimal Initial Improvements:} The initial ResNet18 model showed minimal improvement over the basic CNN model, suggesting limitations due to training, dataset, or model configuration.
    \item Time pressures caused by coursework commitments and personal circumstances resulting in reduced project progress during later weeks.
\end{itemize}

\subsection*{Risks}\label{risks}
\begin{itemize}
    \tightlist
    \item \textbf{Scope of Architectures and Techniques:} There are many different architectures and techniques that could be explored, which could risk spreading efforts too thin. \textbf{Mitigation:} Will narrow the focus to a select few approaches during the beginning of Semester 2, allowing deeper experimentation and more detailed analysis.
    \item \textbf{Evaluation Criteria:} Uncertainty in knowing when a model is performing satisfactorily or when further experimentation is worthwhile. \textbf{Mitigation:} Use clear performance metric thresholds (e.g. recall, accuracy, F1 score) and monitor improvement trends to decide when to stop.
    \item \textbf{Overfitting:} The relatively small dataset increases the risk of models overfitting to the training data, particularly when using high-capacity deep learning models. \textbf{Mitigation:} Will confirm suitable cut-off epochs by plotting and comparing both training and validation loss for convergence; will also experiment with k-fold cross-validation to reduce bias.
    \item \textbf{Computational Resources:} Larger, potentially more effective models (e.g., deeper ResNet architectures or ensemble models) require significant GPU resources, which are currently a limitation. \textbf{Mitigation:} Utilise University computing resources (e.g., ST Linux computers with GPUs or SoCS clusters) to enable the training and evaluation of large models.
    \item \textbf{Workload Balancing:} Balancing the intensive model experimentation and write-up required for the dissertation with coursework from other classes, especially towards the end of the semester. \textbf{Mitigation:} Establishing a clear timeline for the second semester and the holidays, prioritising dissertation work to ensure timely completion.
\end{itemize}
    
\section*{Plan}\label{plan}
\begin{itemize}
    \tightlist
    \item 
      \textbf{Winter Break:} Test and implement a range of seed values and training:validation:test splits; refine current models (ResNet18, CBAM, NCA/kNN); implement ensemble model; confirm overfitting and suitable cut-off epochs by plotting and analysing training/validation loss; create a comparison notebook; continue working on background section write-up. \textbf{Deliverables:} Updated models, ensemble model, comparison notebook, and initial background section write-up.
    \item
      \textbf{Weeks 1-2:} Explore additional model architectures and experiments (e.g., Inception, deeper ResNet variants); experiment with segmentation approaches if progress allows; continue working on write-up. \textbf{Deliverables:} Implemented additional models, initial segmentation experiments (if feasible), updated comparison notebook.
    \item
      \textbf{Weeks 3-4:} Continue experimenting with additional architectures and refining promising models; continue working on write-up. \textbf{Deliverables:} Updated models, refined configurations, expanded comparison notebook, completion of background section.
    \item
      \textbf{Weeks 5-6:} Refine selected models based on performance metrics; perform deeper analysis to identify optimal configurations; continue working on write-up. \textbf{Deliverables:} Refined models, updated performance metrics, detailed comparisons.
    \item 
      \textbf{Week 7:} Conduct any final experiments, confirm results, and complete outstanding analyses; continue working on write-up. \textbf{Deliverables:} Finalised experimental results and validated evaluations.
    \item
      \textbf{Weeks 8-10:} Dissertation write-up only. \textbf{Deliverable:} First draft submitted to supervisor two weeks before final deadline.
\end{itemize}



\section*{Ethics}
This project uses an openly available dataset under a Creative Commons (CC BY-NC-SA 4.0) license, 
permitting non-commercial use, adaptation, and redistribution with appropriate credit and 
share-alike conditions. The dataset contains no personally identifiable information, and all 
experiments comply with the licensing terms. No additional ethical approval is required.

\end{document}
